\section{Spectral Sequences for Higher Genus}

\subsection{The Leray Spectral Sequence for the Universal Curve}
\index{Leray spectral sequence!for universal curve|textbf}
\index{Gauss-Manin connection!in Leray spectral sequence}

For the universal curve $\pi: \mathcal{C}_g \to \mathcal{M}_g$, the Leray spectral sequence reads:
$$E_2^{p,q} = H^p\bigl(\mathcal{M}_g,\, R^q\pi_*\mathbb{C}\bigr) \;\Longrightarrow\; H^{p+q}(\mathcal{C}_g,\, \mathbb{C})$$
The local systems $R^q\pi_*\mathbb{C}$ carry the Gauss--Manin connection, and the differentials encode:
\begin{itemize}
\item $d_2$: Gauss--Manin connection (first obstruction to local-to-global extension)
\item $d_3$: Massey products / Kodaira--Spencer class
\item $d_r$ ($r \geq 4$): Higher deformation data
\end{itemize}

\subsection{The Bar Complex Spectral Sequence}

$E_2^{p,q} = H^p(\overline{\mathcal{M}}_{g,n}, \underline{H}^q(\bar{B}^{(g)}(\mathcal{A}))) \Rightarrow H^{p+q}(\bar{B}^{\text{total}}(\mathcal{A}))$

where $\underline{H}^q$ denotes the local system of bar cohomology groups.

\subsection{Convergence and Degeneration}

\begin{theorem}[Convergence Criterion; \ClaimStatusProvedHere]\label{thm:convergence-criterion-spectral}
\index{spectral sequence!bar complex convergence|textbf}
The spectral sequence
\[
E_2^{p,q} = H^p(\overline{\mathcal{M}}_{g,n}, \underline{H}^q(\bar{B}^{(g)}(\mathcal{A}))) \Rightarrow H^{p+q}(\bar{B}^{\mathrm{total}}(\mathcal{A}))
\]
converges if:
\begin{enumerate}
\item The chiral algebra $\mathcal{A}$ is rational (finitely many irreducible modules)
\item The moduli space $\overline{\mathcal{M}}_{g,n}$ is replaced by its Deligne-Mumford compactification
\end{enumerate}
These conditions ensure convergence of the spectral sequence at each \emph{fixed} genus $g$ (by bounded filtration, see Step~2 below).  The convergence of the full genus sum $\sum_{g \geq 0} g_s^{2g-2} H^*(\bar{B}^{(g)})$ is a separate issue requiring $|g_s| < \epsilon(\mathcal{A})$ for a radius of convergence $\epsilon(\mathcal{A})$ determined by the growth of $\dim H^*(\bar{B}^{(g)})$ with $g$.
\end{theorem}

\begin{proof}
We verify convergence using the bounded filtration criterion for spectral sequences of sheaf cohomology on proper varieties.

\textbf{Step 1: Finite-dimensionality of $E_2$.}
Condition (1) (rationality) implies that $\mathcal{A}$ has finitely many irreducible modules $\{V_1, \ldots, V_N\}$.  The bar complex $\bar{B}^{(g)}(\mathcal{A})$ decomposes as a direct sum over fusion channels: $\bar{B}^{(g)}_n(\mathcal{A}) = \bigoplus_{\vec{i}} \bar{B}^{(g)}_n(\mathcal{A})_{\vec{i}}$ where $\vec{i} = (i_1, \ldots, i_{n+1}) \in \{1, \ldots, N\}^{n+1}$.  The local system $\underline{H}^q(\bar{B}^{(g)}(\mathcal{A}))$ on $\overline{\mathcal{M}}_{g,n}$ therefore has finite-dimensional stalks (bounded by $N^{n+1} \cdot \dim \bar{B}^{(g)}_n$ in each fusion channel).

Condition (2) (DM compactification) ensures that $\overline{\mathcal{M}}_{g,n}$ is a proper Deligne-Mumford stack of finite type.  By properness, $H^p(\overline{\mathcal{M}}_{g,n}, \mathcal{L})$ is finite-dimensional for any constructible local system $\mathcal{L}$.  Therefore each $E_2^{p,q}$ is finite-dimensional.

\textbf{Step 2: Bounded filtration.}
The $E_2$ page is bounded: $E_2^{p,q} = 0$ for $p > \dim_{\mathbb{R}}(\overline{\mathcal{M}}_{g,n}) = 6g - 6 + 2n$ (by dimension of the moduli space) and for $p < 0$.  In the $q$-direction, rationality gives $\underline{H}^q = 0$ for $q$ outside a finite range (the bar complex at each genus has bounded cohomological amplitude, determined by the conformal weights of the finitely many irreducibles).  A spectral sequence with $E_2$ page bounded in a finite rectangle necessarily converges: all differentials $d_r$ for $r$ exceeding the width of the rectangle are zero.

\textbf{Step 3: Genus summation.}
The total bar complex sums over genera: $\bar{B}^{\mathrm{total}}(\mathcal{A}) = \bigoplus_{g \geq 0} g_s^{2g-2} \bar{B}^{(g)}(\mathcal{A})$.  Condition (2) ($|g_s| < \epsilon(\mathcal{A})$) ensures this sum converges.  The bound $\epsilon(\mathcal{A})$ is determined by the growth rate of $\dim H^*(\overline{\mathcal{M}}_{g,n}, \underline{H}^q(\bar{B}^{(g)}))$ as a function of $g$.  For a rational algebra, the factorization structure gives the bound $\dim E_2^{p,q}|_{\text{genus } g} \leq C \cdot N^{2g}$ (from the $N$ fusion channels on a genus-$g$ surface with $2g$ independent cycles), so $\epsilon(\mathcal{A}) = N^{-1}$ suffices for absolute convergence of the genus sum.
\end{proof}

\begin{theorem}[Degeneration at $E_2$; \ClaimStatusProvedElsewhere]
For special values of central charge:
\begin{itemize}
\item $c = 0$: Topological theory, degenerates at $E_1$
\item $c = 26$: Critical bosonic string, degenerates at $E_2$  
\item $c = 15$: Critical superstring (matter sector), conjectured degeneration at $E_2$
\end{itemize}
\end{theorem}

\subsection{Connection to Genus Expansion and Feynman Diagrams}

\begin{remark}[Spectral Sequence and Genus Expansion --- Heuristic Interpretation]\label{rem:ss-genus}
The spectral sequence computing $H^*(\bar{B}(\mathcal{A}))$ has pages that heuristically correspond to genus contributions:

\textbf{Page $E_r$}: Contributions from graphs/diagrams with up to $(r-1)$ loops.

The heuristic correspondence (to be made precise for specific chiral algebras) is:
\begin{itemize}
\item $E_1^{p,q}$: Tree-level contributions (genus 0)
\item $E_2^{p,q}$: Includes one-loop corrections (genus 1)
\item $E_3^{p,q}$: Includes two-loop corrections (genus 2)
\item $E_\infty^{p,q}$: All genera
\end{itemize}

\textbf{Physical interpretation}:
$$E_r \approx \text{Feynman graphs with } \leq (r-1) \text{ loops}$$

The differentials $d_r$ implement quantum corrections by integrating over moduli spaces of curves:
$$d_r: E_r^{p,q} \to E_r^{p+r,q-r+1}$$
corresponds to
$$\int_{\overline{\mathcal{M}}_{g,n}} \omega_{\text{correlator}}$$
where $g = r-1$ is the genus.

The filtration on $\bar{B}(\mathcal{A})$ by pole order along collision divisors corresponds to the loop expansion: order-$0$ pole = tree diagram; order-$n$ pole = $n$ internal loops.  The $E_r$ page consists of equivalence classes of diagrams modulo those with $< r$ loops.
\end{remark}

\begin{remark}[Scope and relation to established results]\label{rem:ss-genus-scope}
The page-by-page genus correspondence stated above is a heuristic physical interpretation, not a theorem.  What has been proved rigorously in this monograph is the following:
\begin{enumerate}
\item The Prism Principle (Corollary~\ref{cor:prism-principle}) establishes that the bar complex decomposes a chiral algebra into its operadic spectrum, with genus-0 structure governed by $\mathrm{HyCom} \leftrightarrow \mathrm{Grav}$ Koszul duality and higher-genus structure by the Feynman transform of the modular operad $\{\overline{\mathcal{M}}_{g,n}\}$ (Theorem~\ref{thm:prism-higher-genus}).
\item The genus filtration (Theorem~\ref{thm:genus-induction-strict}) ensures $d^2 = 0$ at all genera via the modular operad axioms.
\item The Koszul complementarity theorem (Theorem~\ref{thm:koszul-complementarity-higher-genus}) gives the precise decomposition $Q_g(\mathcal{A}) \oplus Q_g(\mathcal{A}^!) \cong H^*(\overline{\mathcal{M}}_g, Z(\mathcal{A}))$ of quantum corrections at each genus.
\end{enumerate}
What remains conjectural is the \emph{identification of the spectral sequence page $E_r$ with genus-$(r-1)$ contributions specifically}.  The rigorous Feynman transform framework organizes all genera simultaneously rather than page-by-page; making the page-genus dictionary precise would require an explicit comparison between the pole-order filtration and the genus grading in the Feynman transform, which has not been carried out in general.
\end{remark}

\subsection{Computational Tools}

The differentials can be computed via:
\begin{enumerate}
\item \textbf{Čech cohomology:} Cover $\overline{\mathcal{M}}_{g,n}$ by affine opens
\item \textbf{Dolbeault cohomology:} Use $\bar{\partial}$-operator techniques
\item \textbf{Combinatorial models:} Jenkins-Strebel differentials
\item \textbf{Topological recursion:} Eynard-Orantin formalism
\end{enumerate}

\subsection{Spectral Sequence for Bar Complex}

\begin{theorem}[Bar Spectral Sequence; \ClaimStatusProvedElsewhere]
The filtration by configuration degree yields a spectral sequence:
$$E_1^{p,q} = H^q(\overline{C}_{p+1}(X), j_*j^*\mathcal{A}^{\boxtimes(p+1)}) \Rightarrow H^{p+q}(\bar{B}^{\text{ch}}(\mathcal{A}))$$

\textbf{Key Properties:}
\begin{enumerate}
\item $E_2$ page: Computed by residues at boundary divisors
\item Convergence: Always for finite-type chiral algebras
\item Degeneration: At $E_2$ for Koszul algebras (quadratic with no higher relations)
\item Differential $d_r$: Encodes $(r+1)$-fold collisions
\end{enumerate}

\textbf{Application to Free Fermions:}
\begin{itemize}
\item $E_1^{p,0} = \wedge^p(\mathcal{F} \otimes H^0(X, \omega_X))$ 
\item $d_1 = 0$ (no relations beyond anticommutativity)
\item Collapses at $E_1 = E_{\infty}$
\item Recovers $\bar{B}^{\text{ch}}(\mathcal{F}) = \wedge^{\bullet}(\mathcal{F}[1])$
\end{itemize}

\textbf{Application to W-algebras:}
For $\mathcal{W}_k(\mathfrak{g}, f)$ at admissible level:
\begin{itemize}
\item $E_1$: Free generators from W-currents
\item $E_2$: Normal ordered products and null fields
\item $E_3$: Quantum corrections from BRST cohomology
\item Convergence requires careful analysis of Virasoro representations
\end{itemize}
\end{theorem}

\begin{example}[Computing $E_2$ Page]
For a chiral algebra with generators $\phi_i$ of conformal weight $h_i$:

$$E_2^{p,q} = \frac{\text{Ker}(d_1: E_1^{p,q} \to E_1^{p+1,q})}{\text{Im}(d_1: E_1^{p-1,q} \to E_1^{p,q})}$$

where $d_1$ is computed from OPE residues:
$$d_1(\phi_{i_1} \otimes \cdots \otimes \phi_{i_p}) = \sum_{j<k} \sum_\ell C_{i_j i_k}^\ell \phi_{i_1} \otimes \cdots \widehat{i_j} \cdots \widehat{i_k} \cdots \otimes \phi_\ell$$
\end{example}

\begin{remark}[Physical Interpretation]
In string theory:
\begin{itemize}
\item $E_1$: Off-shell string states
\item $d_1$: BRST operator
\item $E_2$: Physical (on-shell) states
\item Higher pages: Quantum corrections and anomalies
\end{itemize}
\end{remark}
